%!TEX root = ../nauchka.tex

\section{Распараллеливание верификации}

\begin{frame}
\frametitle{Текущее состояние: Data Parallelism}

\begin{block}{AlphaBetaCROWN использует параллелизм по данным}
Алгоритм Branch-and-Bound разбивает входное пространство на подзадачи, каждая решается на отдельном GPU
\end{block}

\vspace{0.5cm}

\textbf{Ограничение:}
\begin{itemize}
    \item Решает проблему \textbf{времени вычислений}
    \item НЕ решает проблему \textbf{памяти}
    \item Если одна задача не влезает в память одного GPU --- верификация невозможна
\end{itemize}

\vspace{0.3cm}

\begin{alertblock}{Необходимость модельного параллелизма}
Нужно разделить саму модель и матрицы $A$ между GPU
\end{alertblock}

\end{frame}

\begin{frame}
\frametitle{Опыт из обучения моделей}

\textbf{Три основных подхода к распределенному обучению:}

\vspace{0.3cm}

\begin{enumerate}
    \item \textbf{Tensor Parallelism (TP)}
    \begin{itemize}
        \item Разделение матриц внутри слоя
        \item Megatron-LM, широко используется для LLM
    \end{itemize}

    \item \textbf{Pipeline Parallelism (PP)}
    \begin{itemize}
        \item Разделение слоев модели по глубине
        \item GPipe, используется для очень глубоких сетей
    \end{itemize}

    \item \textbf{Fully Sharded Data Parallelism (FSDP)}
    \begin{itemize}
        \item Шардирование параметров и оптимизатора
        \item FSDP (PyTorch), ZeRO (DeepSpeed)
    \end{itemize}
\end{enumerate}

\vspace{0.3cm}

\textbf{Вопрос:} Какой подход применим к верификации?

\end{frame}

\begin{frame}
\frametitle{Математика обратного распространения границ}

\begin{block}{Ключевая операция CROWN}
$$A_{prev} = A_{curr} \cdot W$$
где $A$ --- матрица символьных границ, $W$ --- веса слоя
\end{block}

\vspace{0.3cm}

\textbf{Для линейного слоя:}
\begin{itemize}
    \item $A_{curr}$: размер $(B \times S \times N_{out})$
    \item $W$: размер $(N_{out} \times N_{in})$
    \item $A_{prev}$: размер $(B \times S \times N_{in})$
\end{itemize}

\vspace{0.3cm}

\begin{exampleblock}{Аналогия с обучением}
Это похоже на обратное распространение, но вместо градиентов мы передаем \textbf{символьные матрицы границ}
\end{exampleblock}

\end{frame}
